\documentclass[10pt,a4paper]{amsart}
\usepackage[margin=19mm]{geometry}
\usepackage{amsmath,amssymb,mathtools}
\usepackage[scr=rsfs,cal=euler]{mathalfa}
\usepackage{xfrac}
\usepackage[unicode,hidelinks,bookmarks=false]{hyperref}
\newcommand{\ie}{i.e.\ }
\newcommand{\cf}{cf.\ }
\newcommand{\resp}{respectively\ }
\newcommand{\Subsection}{\subsection}
\providecommand{\nobreakdash}{\nobreak\hbox{-}}
\setlength{\emergencystretch}{2em}
\multlinegap0pt
\usepackage{bm}
\usepackage{bbm}
\usepackage{dsfont}
\usepackage{xspace}
\usepackage{cancel}
\usepackage{mathtools}
\usepackage[shortlabels]{enumitem}
\usepackage{amsthm}
\makeatletter
\@ifundefined{thm}{\newtheorem{thm}{Theorem}}{}
\@ifundefined{lemma}{\newtheorem{lemma}[thm]{Lemma}}{}
\@ifundefined{prop}{\newtheorem{prop}[thm]{Proposition}}{}
\makeatother
\title[Fourier characterizations and non-triviality of Gelfand-Shil]{Fourier characterizations and non-triviality of Gelfand-Shilov spaces, with applications to Toeplitz operators}
\author{Albin Petersson}
\date{Source: arXiv:2202.00938v1 (CC BY 4.0)}
\begin{document}
\maketitle

\noindent\emph{Source and licence.} Fourier characterizations and non-triviality of Gelfand-Shilov spaces, with applications to Toeplitz operators. Version arXiv:2202.00938v1 (CC BY 4.0), distributed under the Creative Commons Attribution 4.0 International licence, as recorded in the arXiv licence metadata for this exact version and shown on the abstract page https://arxiv.org/abs/2202.00938v1. Full rights notice: licenses/arxiv-2202.00938-CC-BY-4.0.txt.\par\smallskip

\noindent\textit{Editorial scope.} Counted unit: the Thm whose source label is \texttt{thm:SsSTFTdual} in the article (source0.tex lines 726--740, source label \texttt{thm:SsSTFTdual}), delivered together with the article's own complete original proof of it (source0.tex lines 741--793, one \texttt{proof} environment of 53 lines). No mathematics is written, repaired or re-derived: statement and proof are the article's own text line for line. Required context retained uncounted: prop:Ssalt (lines 135--150); prop:Ssdual (lines 266--280); thm:SsSTFT (lines 397--411); lem:absolutineq (lines 425--431). Every definition, display and label of those lines is retained unchanged. Omitted from this package and not redistributed: 1--134, 151--265, 281--396, 412--424, 432--725, 794--1027. Nothing inside a retained excerpt is omitted or altered: every retained body is delivered exactly as the article's lines are written, so no omission receipt of any kind accompanies this package. Citations: the retained excerpts carry no citation command at all, so no bibliography entry of the article is transcribed and no reference is redistributed. Reference adaptation: none was needed and none was made; every reference inside the delivered unit resolves inside it, so no unresolved reference or citation is produced. Printed slips kept literal: no printed word, symbol, hypothesis or number of the counted statement or of its proof was altered. Layout and class: the article's own class is \texttt{amsart} with packages including inputenc, amsmath, amsthm, mathrsfs, amssymb, enumitem, cite, mathtools, subfiles, xcolor, comment, pdfpages, hyperref; the wrapper is the lane's pinned \texttt{amsart} editorial wrapper at 10pt on A4, which restates the theorem-like environments the retained text needs and adds the article's own macro definitions that the retained text needs (none), with extra wrapper packages (bm, bbm, dsfont, xspace, cancel, mathtools, enumitem). The delivered numbering is the unit's own. Assets: no retained span contains a figure environment, a graphics-inclusion command, a PDF-page inclusion, a table or a TikZ picture; no image file is copied into this package, the package declares no asset dependency and no image file is redistributed with it. Licence: version arXiv:2202.00938v1 of this article is distributed under the Creative Commons Attribution 4.0 International licence, as recorded in the arXiv licence metadata for this exact version and shown on the abstract page https://arxiv.org/abs/2202.00938v1. A full rights notice is delivered at licenses/arxiv-2202.00938-CC-BY-4.0.txt. Characters: every retained excerpt is pure ASCII, so the delivered engine, XeLaTeX with its default Latin Modern text font, reports no missing character.
\begin{prop} \label{prop:Ssalt}
Suppose $s>0$ and $f\in C^\infty(\mathbb{R}^n)$. Then
\begin{enumerate}[label=(\alph*)]
    \item $f\in S_s(\mathbb{R}^n)$ if and only if there are constants $C_\beta,r>0$ such that
    \begin{equation}\label{eq:Ssalt}
    |D^\beta f(x)| \leq C_\beta e^{- r |x|^{1/s}}
    \end{equation}
    for all multi-indices $\beta$;
    \item $f\in\Sigma_s(\mathbb{R}^n)$ if and only if for every $r>0$
\begin{equation}\label{eq:Ssigalt}
    |D^\beta f(x)| \leq C_{r,\beta} e^{- r |x|^{1/s}}
\end{equation}
holds for all multi-indices $\beta$, where $C_{r,\beta}>0$ depends only on $r$ and $\beta$.

\end{enumerate}
\end{prop}

\begin{prop}\label{prop:Ssdual}
 Suppose $1\leq p \leq \infty$.
\begin{enumerate}[label=(\roman*)]
    \item $u\in (S_s)'(\mathbb{R}^n)$ if and only if for every $r>0$ there is an $N\geq 0$ such that
    \begin{equation*}
        |\langle u, f \rangle | \leq C_N \sum_{|\alpha|\leq N} || D^\alpha f e^{r|\cdot|^{1/s}} ||_p,
    \end{equation*}
    for any $f\in S_s(\mathbb{R}^n)$.
     \item $u\in (S^\sigma)'(\mathbb{R}^n)$ if and only if for every $r>0$ there is an $N\geq 0$ such that
    \begin{equation*}
        |\langle u, f \rangle | \leq C_N \sum_{|\alpha|\leq N} || D^\alpha \hat{f} e^{r|\cdot|^{1/\sigma}} ||_p,
    \end{equation*}
    for any $f\in S^\sigma(\mathbb{R}^n)$.
\end{enumerate}
\end{prop}

\begin{thm}\label{thm:SsSTFT}
Suppose $s,\sigma> 0$.
\begin{enumerate}[label=(\roman*)]
    \item Let $\phi\in S_s(\mathbb{R}^n)\setminus\{0\}$. Then $f\in S_s(\mathbb{R}^n)$ if and only if there is an $r>0$ such that 
        \begin{equation}\label{eq:STFT3}
            |V_\phi f(x,\xi) | \leq C_N (1+|\xi|^2)^{-N} e^{- r |x|^{1/s}}
        \end{equation}
        for every $N\geq 0$.
     \item Let $\phi\in S^\sigma(\mathbb{R}^n)\setminus\{0\}$. Then $f\in S^\sigma(\mathbb{R}^n)$ if and only if there is an $r>0$ such that
        \begin{equation}\label{eq:STFT4}
            |V_{\phi} f(x,\xi) | \leq C_N (1+|x|^2)^{-N} e^{- r |\xi|^{1/\sigma}}
        \end{equation}
        for every $N\geq 0$.
\end{enumerate}
\end{thm}

\begin{lemma}\label{lem:absolutineq}
If $s>0$ then there is a $C \geq 1$ such that
\begin{equation}\label{eq:absolutineq}
    C^{-1} (|x|^{1/s} + |y|^{1/s} ) \leq |y|^{1/s} + |y-x|^{1/s} \leq C (|x|^{1/s} + |y|^{1/s} ).
\end{equation}
for every $x,y\in\mathbb{R}^n$.
\end{lemma}

\begin{thm}\label{thm:SsSTFTdual}
% ------------------------------------
%   Thm: (S_s)' STFT
% ------------------------------------
\begin{enumerate}[label=(\roman*)]
    \item Let $\phi\in S_s(\mathbb{R}^n)\setminus\{0\}$ and $f\in \mathscr{D}'(\mathbb{R}^n)$. Then $f\in(S_s)'(\mathbb{R}^n)$ if and only if for every $r>0$ there is an $N_0\geq 0$ such that
    \begin{equation}\label{eq:STFTdual1}
        |V_\phi f(x,\xi)| \leq C_{r}(1+|\xi|^2)^{N_0} e^{r|x|^{1/s}}.
    \end{equation}
     \item Let $\phi\in S^\sigma(\mathbb{R}^n)\setminus\{0\}$ and $f\in \mathscr{F}\mathscr{D}'(\mathbb{R}^n)$. Then $f\in(S^\sigma)'(\mathbb{R}^n)$ if and only if for every $r>0$ there is an $N_0\geq 0$ such that
    \begin{equation}\label{eq:STFTdual2}
        |V_\phi f(x,\xi)| \leq C_{r}(1+|x|^2)^{N_0} e^{r|\xi|^{1/\sigma}}.
    \end{equation}
\end{enumerate}
\end{thm}

\begin{proof}
Suppose $f\in (S_s)'$, $\phi\in S_s\setminus\{0\}$. Then by Proposition \ref{prop:Ssdual} and Leibniz formula, there is an $N_0>0$ such that
\begin{align*}
    |V_\phi f(x,\xi)| &= |(f,\phi(\cdot-x)e^{i\langle\cdot,\xi\rangle})|  \\
    &\leq C_{N_0} \sum_{|\alpha|\leq N_0} || D^\alpha \left(\phi(\cdot-x)e^{i\langle \cdot,\xi\rangle}\right) e^{r|\cdot|^{1/s}}||_2\\
    &= C_{N_0} \sum_{|\alpha|\leq N_0}\sum_{\gamma\leq\alpha}\binom{\alpha}{\gamma} || (D^\gamma \phi)(\cdot-x) \xi^{\alpha-\gamma} e^{i\langle \cdot,\xi\rangle} e^{r|\cdot|^{1/s}}||_2 \\
    &\leq C'_{N_0} (1+|\xi|^2)^{N_0} \sum_{|\alpha|\leq N_0}\sum_{\gamma\leq\alpha}\binom{\alpha}{\gamma} || (D^\gamma \phi)(\cdot-x) e^{r|\cdot|^{1/s}}||_2
\end{align*}
By Theorem \ref{prop:Ssalt}, there is an $r_0>0$ such that $|D^{\gamma}\phi(y-x)|\leq C_\gamma e^{-r_0|y-x|^{1/s}}$, hence
\begin{align*}
    |V_\phi f(x,\xi)| &\leq C'_{N_0} (1+|\xi|^2)^{N_0} \sum_{|\alpha|\leq N_0}\sum_{\gamma\leq\alpha}\binom{\alpha}{\gamma} || e^{-r_0|\cdot-x|^{1/s}} e^{r|\cdot|^{1/s}}||_2.
\end{align*}
By Lemma \ref{lem:absolutineq}, there is a $c\geq 1$ such that
\begin{equation*}
    - r_0|y-x|^{1/s} \leq   r_0 |x|^{1/s} - r_0/c \cdot |y|^{1/s}.
\end{equation*}
Let $r\in (0, r_0 /(2c))$. Then
\begin{equation*}
    - r_0|y-x|^{1/s} \leq - 2 c r|y-x|^{1/s} \leq  2 c r |x|^{1/s} - 2 r |y|^{1/s},
\end{equation*}
and
\begin{align*}
    |V_\phi f(x,\xi)| &\leq C'_{N_0} (1+|\xi|^2)^{N_0} \sum_{|\alpha|\leq N_0}\sum_{\gamma\leq\alpha}\binom{\alpha}{\gamma} || e^{2cr |x|^{1/s} - 2 r |\cdot|^{1/s} + r|\cdot|^{1/s}}||_2 \\
    &=C'_{N_0} (1+|\xi|^2)^{N_0} e^{2 c r |x|^{1/s}} \sum_{|\alpha|\leq N_0}\sum_{\gamma\leq\alpha}\binom{\alpha}{\gamma} ||  e^{-r|\cdot|^{1/s}}||_2  \\
    &\leq C''_{N_0} (1+|\xi|^2)^{N_0} e^{2 c r |x|^{1/s}}.
\end{align*}
Clearly, the inequality still holds if we let $r>r_0/(2c)$ (the right hand side only becomes larger), hence we have shown that the inequality is valid for all $r>0$, as was to be shown.

Now suppose that for every $r>0$ there is an $N_0 \geq 0$ such that \eqref{eq:STFTdual1} holds. Then by Moyal's identity, for every $\varphi \in S_s $
\begin{equation*}
    |(f,\varphi)|\leq ||\phi||_2^{-2} |(V_\phi f, V_\phi \varphi)|, 
\end{equation*}
hence
\begin{equation*}
    |(f,\varphi)|\leq ||\phi||_2^{-2} \int \int |V_\phi f(x,\xi)|\cdot |V_\phi \varphi(x,\xi)| \, d x \, d \xi.
\end{equation*}
By assumption combined with Theorem \ref{thm:SsSTFT}, for every $r,r_1>0$ and any $N_1\geq 0$ there is an $N_0\geq 0$ such that
\begin{align*}
    |(f,\varphi)|&\leq C_{N_0}||\phi||_2^{-2} \int \int  (1+|\xi|^2)^{N_0} e^{r|x|^{1/s}}  |V_\phi\varphi(x,\xi)| \, d x \, d \xi \\
    &= C_{N_0}||\phi||_2^{-2} \int \int  (1+|\xi|^2)^{(N_0-N_1)} e^{(r-r_1)|x|^{1/s}}  \left|V_\phi\varphi(x,\xi)(1+|\xi|^2)^{N_1} e^{r_1|x|^{1/s}}\right| \, d x \, d \xi \\
    &\leq C_{N_0}||\phi||_2^{-2} p^{\phi,s}_{r_1,N_1}(\varphi)\int \int  (1+|\xi|^2)^{-(N_1-N_0)} e^{-(r_1-r)|x|^{1/s}} \, d x \, d \xi.
\end{align*}
Pick $N_1$ such that $N_1>N_0$ and pick $r$ such that $r<r_1$. Then we obtain
\begin{align*}
    |(f,\varphi)|&\leq   C_{N_0,N_1,r}||\phi||_2^{-2} p^{\phi,s}_{r_1,N_1}(\varphi)
\end{align*}
for all $r_1>r$. By picking $r>0$ arbitrarily small, we thus obtain
\begin{align*}
    |(f,\varphi)|&\leq   C'_{N_1,r_1}\cdot p^{\phi,s}_{r_1,N_1}(\varphi)
\end{align*}
for all $r_1>0$. This completes the proof of (i). The proof of (ii) is very similar, utilizing the fact that $\phi\in S^\sigma$ is equivalent to $\hat{\phi}\in S_\sigma $, and the fact that $V_\phi f(x,\xi) = (\hat{f},\hat{\phi}(\cdot-\xi)e^{-i\langle \cdot,x \rangle})$.    
    
\end{proof}

\end{document}
